% concatenated from pv-nconst.tex
\documentclass[a4paper, twoside, 10pt]{article}
\usepackage{amsmath,amsthm,amsfonts,latexsym,amscd,amssymb, enumerate}
\usepackage{hyperref, xcolor, soul, xparse}
\numberwithin{equation}{section}
\setlength{\parskip}{1ex plus 0.5ex minus 0.2ex}
\addtolength{\textwidth}{0.6cm}
\addtolength{\textheight}{0.5cm}
\addtolength{\hoffset}{-0.5cm}
\setlength{\oddsidemargin}{1cm} 
\setlength{\evensidemargin}{1cm} 
\pagestyle{myheadings}
\thispagestyle{empty}
\markboth{\small{M. Aqib, H. M. Shah and S. Deshmukh}}
{\small{  }}

\newtheorem{theorem}{Theorem}[section]
\newtheorem{lemma}{Lemma}[section]
\newtheorem{corollary}{Corollary}[section]
\newtheorem{proposition}{Proposition}[section]
\theoremstyle{definition}
\newtheorem{definition}{Definition}[section]
\newtheorem{example}{Example}[section]
\newtheorem{remark}{Remark}[section]
\title{\textbf{Non-compact Ricci Solitons of Finite Volume with Potential Field of Constant Length}}                                 
\author{Hemangi Madhusudan Shah, Sharief Deshmukh, and Mohammad Aqib}            
\date{}

\begin{document}

\maketitle
\begin{abstract}
We study non-compact Ricci solitons of finite volume whose potential
vector field has constant length. Under the assumptions that the scalar 
curvature is constant along the integral curves of the potential field 
and that a natural divergence term is integrable on the unit tangent bundle, we prove that such Ricci solitons are necessarily trivial. As applications, we obtain rigidity and non-existence results for Ricci
solitons whose potential field is the Reeb vector field of almost contact 
metric and almost $\alpha$-cosymplectic manifolds. In dimension three,
we derive consequences for almost $\alpha$-cosymplectic and contact
metric manifolds, and we compare our results with the classification of  
homogeneous almost $\alpha$-cosymplectic Ricci solitons due to Li and Liu.
Several examples and non-examples are included to illustrate the necessity of the finite-volume and sign assumptions.
\end{abstract}

\vskip.1in
 
%\noindent{\sc Corresponding author:  Hemangi Madhusudan %Shah\\ E-mail:	hemangimshah@hri.res.in} 

 \vskip.1in
\noindent {{\sc 2020 AMS Mathematics Subject Classification}: Primary: 53C20; 53C21; 53C25.}  

\vskip.1in
\noindent {\sc Key Words}: Finite volume Ricci solitons; Ricci solitons; Incompressible vector field




\section{Introduction}

Recall a Ricci soliton is the quadruple $\left( M^{n},g,X,\lambda \right) $,
where $\left( M^{n},g\right) $ is an $n$-dimensional Riemannian manifold, $X$
is a smooth vector field, $\lambda$ is a constant, satisfying
\begin{equation}\label{e0}
\frac{1}{2}\mathcal{L} _{X}g+\operatorname{Ric}=\lambda g\text{,}  
\end{equation}
where $\mathcal{L} _{X}g$ is the Lie derivative with respect to $X$ and $\operatorname{Ric}$ is
the Ricci tensor of the Riemannian manifold $\left( M^{n},g\right) $. 
If $X$ is a Killing vector field, then the Ricci soliton $\left(M^{n},g, X,\lambda \right)$ is an Einstein manifold, and in this case it is called a trivial soliton.

This paper aims to study non-compact Ricci solitons of finite volume whose potential vector field has constant length. This condition appears naturally in almost contact geometry, where the Reeb vector field is a unit vector field. Specifically, studying Ricci solitons with a potential field of constant length reveals the geometry of a huge family of almost contact and almost paracontact manifolds, including contact, $K$-contact, Sasakian, Kenmotsu, and trans-Sasakian manifolds. Our first main result shows that, under a natural integrability
condition and a scalar-curvature condition along the flow of the potential field; such solitons are necessarily trivial.

Here, we would like to point out that studying Ricci solitons $\left(
M^{n},g,X,\lambda \right) $ with potential field $X$ of constant length
reveals the geometry of almost contact manifolds as well as almost paracontact manifolds, as Ricci solitons with a Reeb vector field, the potential vector field. Moreover, the family of almost contact manifolds and
almost paracontact manifolds are huge, comprising contact manifolds,
paracontact manifolds, K-contact manifolds, Sasakian and Para-Sasakian
manifolds, Quasi-Sasakian manifolds, Kenmotsu manifolds and trans-Sasakian manifolds.

The main results are as follows:
\begin{itemize}
\item Theorem \ref{main}: finite-volume rigidity for constant-length potential fields.
\item Theorem \ref{thm:monotone}: a monotonicity version assuming only $X(S)\le0$.
\item Corollary \ref{cor:gradient}: rigidity for gradient Ricci solitons.
\item Corollary \ref{c2}: the almost $\alpha$-cosymplectic Reeb-potential case.
\item Theorem \ref{t:contact}: classification of three-dimensional contact Ricci solitons with Reeb potential.
\end{itemize}

The paper is organised as follows. In Section 2, we collect basic
identities for Ricci solitons whose potential field has constant
length, and we recall the necessary background on almost contact and
almost $\alpha$-cosymplectic manifolds. In Section 3, we prove the main
finite-volume rigidity theorems. Section 4 contains applications to
almost contact and almost $\alpha$-cosymplectic geometry. In Section 5,
we study contact Ricci solitons in dimension three. Finally, Section 6 contains examples and non-examples illustrating the sharpness of the hypotheses.


%%%%%%%%%%%%%%%%%%%%%%%%%%%%%%%%%%%%%%%%%%%%%%%%%
\section{Preliminaries and Basic Identities}
%%%%%%%%%%%%%%%%%%%%%%%%%%%%%%%%%%%%%%%%%%%%%5


%----------------------------------------------------------------%
\subsection{Ricci solitons with constant-length potential field}
%----------------------------------------------------------------%
Let $\left(M^{n},g,X, \lambda \right) $ be an $n$-dimensional Ricci soliton.
We shall have
\begin{equation}\label{e1}
{\nabla} _{U}X=\lambda U-QU+\phi U\text{,\quad }U\in \Gamma \left(
TM^{n}\right) \text{,}  
\end{equation}
where $\lambda $ is a constant and $Q$ is the Ricci operator and $\phi$ is a skew-symmetric tensor related to $X$ through its dual $1$-form $\eta$ by
\[
\frac{1}{2}d\eta \left( U,V\right) =g\left( \phi U,V\right) \text{,\quad }%
U,V\in \Gamma \left( TM^{n}\right) \text{.}
\]
Suppose that $X$ has constant length. Then by virtue of equation (\ref{e1})
\[
0=U\left( \left\Vert X\right\Vert ^{2}\right) =2g\left( \lambda U-QU+\phi
U,X\right) \text{,\quad }U\in \Gamma \left( TM^{n}\right) \text{,}
\]
that is,
\begin{equation}\label{e2}
\lambda X-QX -\phi X=0\text{.}  
\end{equation}


Taking the inner product of \eqref{e2} with $X$ and using the
skew-symmetry of $\phi$, we obtain
\[
\operatorname{Ric}(X,X)=\lambda \|X\|^2.
\]
Thus, if $\|X\|^2=c>0$, then
\begin{equation}\label{e3}
\operatorname{Ric}\left( X,X\right) =c\lambda \text{,}  
\end{equation}
which is constant whenever $\|X\|$ is constant. Equation (\ref{e2}) will be one of the main tools in studying the geometry of such Ricci solitons.


\subsection{Fundamental Identities}



Differentiating equation (\ref{e2}), while using equation (\ref{e1}), we see
\begin{eqnarray*}
\lambda \left( \lambda U-QU+\phi U\right)  &=&\left( \nabla Q\right) \left(U,X\right) +Q\left( \lambda U-QU+\phi U\right)  \\ 
&&+\left( \nabla \phi \right) \left( U,X\right) +\phi \left( \lambda
U-QU+\phi U\right) \text{.}
\end{eqnarray*}
Choosing a local orthonormal frame $\left\{ u_{1},..,u_{n}\right\} $ and
taking $U=u_{i}$ and on taking the inner product with $u_{i}$ and summing
the resulting equation, while keeping in mind the symmetry of $Q$ and skew
symmetry of $\phi $, we land at
\begin{equation}\label{e4}
n\lambda ^{2}-\lambda S=\frac{1}{2}X(S)+\lambda S-\left\Vert Q\right\Vert
^{2}-\sum\limits_{i}g\left( X,\left( \nabla \phi \right) \left(
u_{i},u_{i}\right) \right) -\left\Vert \phi \right\Vert ^{2}\text{,}  
\end{equation}
where we have used the following
\[
\sum\limits_{i}g\left( Qu_{i},\phi u_{i}\right) =0\text{,\quad }
\sum\limits_{i}g\left( Qu_{i},Qu_{i}\right) =\left\Vert Q\right\Vert ^{2}
\text{,\quad }\sum\limits_{i}g\left( \phi u_{i},\phi u_{i}\right)
=\left\Vert \phi \right\Vert ^{2}\text{.}
\]
Rearranging equation (\ref{e4}), we get%
\begin{eqnarray}\label{e5}
\left\Vert Q\right\Vert ^{2}-\frac{1}{n}S^{2} &=&-\frac{1}{n}\left(
S-n\lambda \right) ^{2}+\frac{1}{2}X(S)   \\
&&-\sum\limits_{i}g\left( X,\left( \nabla \phi \right) \left(
u_{i},u_{i}\right) \right) -\left\Vert \phi \right\Vert ^{2}\text{.} 
\nonumber
\end{eqnarray}

Now, we use equation (\ref{e1}) in computing $\operatorname{div}\left(\phi X\right)$,
we have
\[
\operatorname{div}\left(\phi X \right) =\sum\limits_{i}g\left( \nabla _{u_{i}}\phi
X,u_{i}\right) =\sum\limits_{i}g\left( \left( \nabla \phi \right) \left(
u_{i},X\right) +\phi \left( \nabla _{u_{i}}X\right) ,u_{i}\right) \text{,} 
\]
that is,
\[
\operatorname{div}\left( \phi X\right) =\sum\limits_{i}g\left( \left( \nabla \phi
\right) \left( u_{i},X\right) +\phi \left( \lambda u_{i}-Qu_{i}+\phi
u_{i}\right) ,u_{i}\right) \text{.} 
\]
Using skew symmetry of $\phi $, above equation reduces to
\begin{eqnarray}\label{e6}
\operatorname{div}\left(\phi X \right) = -\sum \limits_{i} g\left(X,\left(\nabla \phi \right) \left( u_{i},u_{i}\right) \right) -\left\Vert \phi \right\Vert
^{2}\text{,}  
\end{eqnarray}
where we have used
\[
\sum\limits_{i}g\left( Qu_{i},\phi u_{i}\right) =0\text{.} 
\]
Inserting equation (\ref{e6}) in equation (\ref{e5}), we arrive at
\begin{equation}\label{e7}
\left\Vert Q\right\Vert ^{2}-\frac{1}{n}S^{2}=-\frac{1}{n}\left( S-n\lambda
\right) ^{2}+\frac{1}{2}X\left( S\right) +\operatorname{div}\left(\phi X \right) 
\text{.}  
\end{equation}

This identity will be the main tool in the proof of our rigidity
theorems.

\subsection{Almost contact metric manifolds}

An almost contact metric manifold $\left( M^{2n+1},g, \varphi,\xi ,\eta
\right) $ comprises a $(2n+1)$-dimensional Riemannian manifold $\left(
M^{2n+1},g\right)$, equipped with a $(1,1)$ tensor field $\varphi$, a unit
vector field $\xi $ called Reeb vector field and its dual $1$-form $\eta$
satisfying
\[
 {\varphi}^{2}=-I+\eta \otimes \xi \text{,\quad } \varphi \xi =0\text{,\quad }%
\eta \circ \varphi =0\text{,\quad }g\left( \varphi U, \varphi V\right)
=g\left( U,V\right) -\eta \left( U\right) \eta \left( V\right) \text{.}
\]

An almost contact metric manifold $\left( M^{2n+1},g,\varphi ,\xi ,\eta
\right) $ is said to be a $K$-contact manifold if the Reeb vector field $\xi 
$ is Killing. Moreover, an almost contact metric manifold $\left(
M^{2n+1},g,\varphi ,\xi ,\eta \right) $ is said to be a Sasakian manifold if
it satisfies
\[
\left( \nabla \varphi \right) \left( U,V\right) =\eta \left( V\right)
U-g\left( U,V\right) \xi 
\]
and it follows that a Sasakian manifold is a $K$-contact manifold, while the
converse is not true. In contact geometry, one of the interesting questions
is to find conditions under which an almost contact metric manifold is a $K$%
-contact manifold and a similar important question is to find conditions under
which a $K$-contact manifold is a Sasakian manifold. In this respect, in \cite{boyer2001einstein}
authors have proved that a compact Einstein $K$-contact manifold is a
Sasakian Einstein manifold (cf. Theorem A in \cite{boyer2001einstein}).


\subsection{Almost \texorpdfstring{$\alpha$}{alpha}-cosymplectic manifolds}

 An almost contact metric manifold 
$\left(M^{2n+1},g, \varphi, \xi ,\eta \right)$ 
is called an  \textit{almost $\alpha$-cosymplectic manifold}, if $d\eta=0$ and $d \Phi=2\alpha\eta\wedge \Phi$, for some constant $\alpha$, where $2$-form $\Phi$ is defined as $\Phi(U,V)=g(U, \varphi V)$. 
In case $\alpha=0$,  an almost $\alpha$-cosymplectic manifold is \textit{cosymplectic}. 
An $\alpha$-cosymplectic manifold is a {\it normal} almost $\alpha$-cosymplectic manifold. 
For more details on almost cosymplectic and cosymplectic manifolds, we refer to \cite{perrone2018classification}
and references therein.
 We observe that the term 
``almost coK\"{a}hler manifold" is referred to as
``almost cosymplectic manifold", in the article
\cite{perrone2018classification}.
Cosymplectic/almost cosymplectic manifolds are odd-dimensional equivalent concepts of K\"{a}hler/almost K\"{a}hler manifolds. 


\bigskip 


An $\alpha$ cosymplectic $3$-manifold
admits a $\varphi$-basis \cite{blair2002riemannian} as:
\begin{proposition}[p. 70, \cite{perrone2018classification}]
On almost $\alpha$-cosymplectic $3$-manifold
  $M^{3}(\varphi,\xi,\eta,g)$,  there exists an orthonormal $\varphi$-basis $\{e,\varphi e, \xi \}$ satisfying
  \begin{eqnarray*}
  he=\sigma e,~~~~
  h \varphi e=-\sigma \varphi e,~~~~
  h\xi=0,
  \end{eqnarray*}
  with $\sigma$ a local smooth eigenfunction of $h$.
\end{proposition}

We would require the following table of the Levi-Civita connection.

\begin{theorem}[p. 70, \cite{perrone2018classification}]\label{lc} 
The covariant derivatives with respect to $\varphi$ basis of almost $\alpha$-cosymplectic $3$-manifold are given by,
  \begin{equation}\label{eq5}
    \begin{cases}
       & \nabla_ee=-a \varphi e-\alpha\xi,~~~~ \nabla_{\varphi e} e=-b \varphi e+\sigma\xi,~~~~ \nabla_\xi e=\mu \varphi e,  \\
       & \nabla_e \varphi e=ae+\sigma\xi,~~~~ \nabla_{\varphi e}\varphi  e=be-\alpha\xi,~~~~ \nabla_\xi \varphi e=-\mu e,  \\
       & \nabla_e\xi =\alpha e-\sigma \varphi e,~~~~ \nabla_{\varphi e}\xi=-\sigma e+\alpha \varphi e,~~~~ \nabla_\xi\xi=0,
    \end{cases}
  \end{equation}
   where $a=g(\nabla_e \varphi e, e)$, $b=-g(\nabla_{\varphi e}e, \varphi e)$ and $\mu=g(\nabla_\xi e, \varphi e)$ are smooth functions.
\end{theorem}

\begin{proposition}\cite{perrone2018classification} \label{r}
The scalar curvature $r$ of almost $\alpha$-cosymplectic $3$-manifold is expressed by
\begin{equation}
r=-6 \alpha^{2}-\operatorname{tr} h^{2}-2\left(a^{2}+b^{2}\right)-2(\varphi e)(a)+2 e(b). 
\end{equation}
\end{proposition}


%%%%%%%%%%%%%%%%%%%%%%%%%%%%%%%%%%%%%%%%%%%%%%%%%%%$
\section{Finite-Volume Rigidity Theorems}
%%%%%%%%%%%%%%%%%%%%%%%%%%%%%%%%%%%%%%%%%%%%%%%%%%%%%

\subsection{Main finite-volume rigidity theorem}


\begin{theorem}\label{main}
Let $\left( M^{n},g,X,\lambda \right) $ be an $n$-dimensional non-compact
Ricci soliton of finite volume without boundary with length of potential field $X$ a constant. If the scalar curvature $S$ is constant along the
integral curves of $X$ and the function $f=\left\langle \nabla _{v}\phi
X,v\right\rangle $ is integrable on the unit tangent bundle $SM^{n}$ of $
M^{n}$, then $\left( M^{n},g,X,\lambda \right) $ is trivial.
\end{theorem}


\begin{proof}
    Assume that the Ricci soliton $\left(M^{n},g,X,\lambda \right)$ is
non-compact without boundary and has finite volume, and that the function 
$f=\left\langle \nabla _{v}\phi X,v\right \rangle$ is integrable on the unit
tangent bundle $SM^{n}$ of $M^{n}$ and therefore, we have 
$\operatorname{div}\left(\phi X \right)$ is integrable on $M^{n}$ and that
\begin{eqnarray*}
\int_{M^{n}}\operatorname{div} \left(\phi X\right)dV =0.
\end{eqnarray*}
Now, integrating equation (\ref{e7}) and assuming that the scalar curvature $S$
is a constant along the integral curves of the potential field $X$, we
conclude
\begin{equation}\label{e8}
\int\limits_{M^{n}}\left( \left\Vert Q\right\Vert ^{2}-\frac{1}{n}
S^{2}\right) dV=-\frac{1}{n}\int\limits_{M^{n}}\left(S-n\lambda \right)
^{2}dV\text{.}  
\end{equation}
Note that the Cauchy-Schwarz inequality implies
\begin{equation}\label{e9}
\left\Vert Q\right\Vert ^{2}\geq \frac{1}{n}S^{2}  
\end{equation}
and thus, the left hand side in equation (\ref{e8}) is non-negative whereas, the
right hand side, it is non-positive. Hence, we conclude
\begin{equation}\label{e10}
\int \limits_{M^{n}}\left(\left\Vert Q\right\Vert ^{2}-\frac{1}{n} S^{2}\right) dV=0\text{,\quad }\int\limits_{M^{n}}\left( S-n\lambda \right)
^{2}dV=0  
\end{equation}
and it amounts to equality in the inequality (\ref{e9}), which holds, if and only
if $Q=\frac{S}{n}I$ and the second equation in (\ref{e10}) implies $S=n\lambda $.
Thus, we conclude
\[
Q=\lambda I\text{,} 
\]%
that is, $\operatorname{Ric}=\lambda g$, and by defining the equation of the Ricci soliton, we get
\[
\mathcal{L} _{X}g=0 
\]
and consequently that $\left( M^{n},g,X,\lambda \right) $ is a trivial Ricci
soliton.
\end{proof}



\subsection{Monotone scalar curvature theorem}


We now strengthen Theorem \ref{main} by relaxing the condition on the
scalar curvature. Instead of requiring that $S$ be constant along the
integral curves of $X$, it suffices that $S$ be non-increasing along these curves.




\begin{theorem}\label{thm:monotone}
Let \((M^{n},g, X,\lambda)\) be an \(n\)-dimensional non-compact Ricci
soliton of finite volume and without boundary whose potential field \(X\) has positive constant length. If the scalar curvature $S$ is non-increasing along the integral curves of $X$,
that is,
\[
X(S)\leq 0,
\]
and the function
$f=\left\langle \nabla _{v}\phi X,v\right\rangle$ is integrable on the
unit tangent bundle $SM^{n}$ of $M^{n}$, then
$\left( M^{n},g,X,\lambda \right)$ is trivial. Moreover,
$S=n\lambda$ and $X(S)=0$ everywhere on $M^{n}$.
\end{theorem}

\begin{proof}
Proceeding as in the proof of Theorem \ref{main}, the hypotheses that
$\left\Vert X\right\Vert$ is constant, that $M^{n}$ is non-compact of
finite volume without boundary, and that $f$ is integrable on
$SM^{n}$ imply
\[
\int_{M^{n}}\operatorname{div}\left(\phi X\right)dV=0.
\]
Integrating the fundamental identity \eqref{e7} over $M^{n}$, we
obtain
\begin{equation}\label{mono_eq}
\int_{M^{n}}\!\left(\left\Vert Q\right\Vert ^{2}
-\frac{1}{n}S^{2}\right)dV
\;=\;
-\frac{1}{n}\int_{M^{n}}\!(S-n\lambda)^{2}\,dV
\;+\;\frac{1}{2}\int_{M^{n}}\!X(S)\,dV.
\end{equation}
By the Cauchy--Schwarz inequality,
$\left\Vert Q\right\Vert^{2}\geq\frac{1}{n}S^{2}$ everywhere on
$M^{n}$, so the left-hand side of \eqref{mono_eq} is non-negative. Since
\[
X(S)
=
2\left(\|Q\|^2-\frac1nS^2\right)
+\frac{2}{n}(S-n\lambda)^2
-2\operatorname{div}(\phi X),
\]
equation \eqref{e7} shows that \(X(S)\) is integrable on \(M^n\).
On the right-hand side, the first term is non-positive and, by
hypothesis, the second term satisfies
\[
\frac{1}{2}\int_{M^{n}}X(S)\,dV\;\leq\;0.
\]
Therefore, the right-hand side of \eqref{mono_eq} is non-positive.
Since the left-hand side of \eqref{mono_eq} is non-negative and the
right-hand side is non-positive, both sides of \eqref{mono_eq} must
vanish. Hence,
\[
\int_{M^n}\left(\|Q\|^2-\frac1nS^2\right)dV=0
\]
and
\[
-\frac1n\int_{M^n}(S-n\lambda)^2\,dV
+\frac12\int_{M^n}X(S)\,dV=0.
\]
The two terms in the latter identity are both non-positive;
therefore each must vanish separately.
\begin{equation}\label{mono_vanish}
\int_{M^{n}}\!\left(\left\Vert Q\right\Vert ^{2}
-\frac{1}{n}S^{2}\right)dV=0,\qquad
\int_{M^{n}}\!(S-n\lambda)^{2}\,dV=0,\qquad
\int_{M^{n}}\!X(S)\,dV=0.
\end{equation}
The first identity, combined with
$\left\Vert Q\right\Vert^{2}\geq\frac{1}{n}S^{2}$, implies
$\left\Vert Q\right\Vert^{2}=\frac{1}{n}S^{2}$ pointwise; this
holds if and only if $Q=\frac{S}{n}I$. The second identity implies
$S=n\lambda$ everywhere. Together, $Q=\lambda I$, that is,
$\operatorname{Ric}=\lambda g$.

Substituting $\operatorname{Ric}=\lambda g$ into the soliton equation
\eqref{e0} yields $\mathcal{L}_{X}g=0$, so $X$ is a Killing vector
field. Consequently, $\left( M^{n},g,X,\lambda \right)$ is a trivial
Ricci soliton.

Finally, \(X(S)\) is continuous and non-positive.
If \(X(S)(p)<0\) at some point \(p\in M^n\), then
\(X(S)<0\) on a neighbourhood of \(p\), which would imply
\[
\int_{M^n}X(S)\,dV<0,
\]
contrary to \eqref{mono_vanish}. Hence \(X(S)=0\) everywhere.
\end{proof}

\bigskip


\begin{remark}
Theorem \ref{thm:monotone} properly extends Theorem \ref{main}. Indeed,
the assumption \(X(S)=0\) is the special case of the weaker hypothesis
\(X(S)\le0\). The proof shows that under the remaining assumptions,
the inequality is rigid: necessarily \(X(S)\equiv0\). Thus every
soliton satisfying \(X(S)\le0\) automatically satisfies the hypothesis
of Theorem \ref{main}.
\end{remark}



\subsection{Integral version}

\begin{proposition}\label{prop:integral-monotone}
Under the assumptions of Theorem \ref{thm:monotone}, the conclusion
remains valid if the pointwise assumption \(X(S)\le0\) is replaced by
the weaker requirements that \(X(S)\in L^1(M)\) and
\[
\int_{M^n}X(S)\,dV\le0.
\]
\end{proposition}


\begin{proof}
The proof is identical to that of Theorem
\ref{thm:monotone}; one only uses the integral inequality
\[
\int_{M^n}X(S)\,dV\le0,
\]
rather than the pointwise condition \(X(S)\le0\).
\end{proof}

As an application, we obtain a clean rigidity result for gradient Ricci
solitons. Recall that a Ricci soliton is called a gradient Ricci
soliton if the potential field is the gradient of a smooth function,
$X=\nabla f$, so that the soliton equation becomes
$\operatorname{Hess}f+\operatorname{Ric}=\lambda g$. Since the dual
$1$-form $\eta=df$ satisfies $d\eta=0$, the skew-symmetric tensor
$\phi$ of equation \eqref{e1} vanishes identically. In particular,
$\operatorname{div}(\phi X)=0$ and the integrability condition on
$SM^{n}$ in Theorem \ref{thm:monotone} is automatically satisfied.

\begin{corollary}\label{cor:gradient}
Let $(M^{n},g,\nabla f,\lambda)$ be a complete, non-compact gradient
Ricci soliton of finite volume without boundary whose gradient \(\nabla f\) has positive constant length. If
\[
\nabla f(S)\leq 0,
\]
then $\nabla f=0$, $f$ is constant, and $(M^{n},g)$ is Einstein with
$\operatorname{Ric}=\lambda g$. In particular, the gradient Ricci
soliton is trivial.
\end{corollary}

\begin{proof}
Since \(X=\nabla f\), the dual $1$-form is
\(\eta=df\), hence \(d\eta=d(df)=0\). Therefore
\(\phi=0\), and the integrability condition in
Theorem \ref{thm:monotone} is automatically satisfied. Hence
Theorem \ref{thm:monotone} yields
\[
\operatorname{Ric}=\lambda g,
\qquad
\mathcal{L}_{\nabla f}g=0.
\]

Because $\nabla f$ is a gradient vector field, its Lie derivative
$\mathcal{L}_{\nabla f}g=2\operatorname{Hess}f$. Hence
$\operatorname{Hess}f=0$, which means that $\nabla f$ is a parallel
vector field.

Suppose, to the contrary, that $\nabla f$ is not identically zero.
Since $\nabla f$ is parallel, it has constant positive length and
defines a nowhere-vanishing parallel vector field on $M^{n}$. By the
de Rham decomposition theorem, $M^{n}$ splits isometrically as
\[
M^{n}\cong N^{n-1}\times\mathbb{R}.
\]
Consequently,
\[
\operatorname{Vol}(M^{n})
=\operatorname{Vol}(N^{n-1})\cdot\operatorname{Vol}(\mathbb{R})
=\infty,
\]
contradicting the finite-volume hypothesis. Therefore
$\nabla f=0$, $f$ is constant, and the gradient Ricci soliton is
trivial.
\end{proof}

\begin{remark}
For a gradient Ricci soliton, the condition
\(\nabla f(S)\le0\) means that the scalar curvature is
non-increasing along the gradient flow lines of the
potential function \(f\). Thus Theorem
\ref{thm:monotone} shows that even this weak monotonicity
assumption forces the soliton to be Einstein and, under the finite-volume hypothesis, trivial.
\end{remark}

\bigskip


We also record the following consequence of Theorem \ref{main}:


\subsection{Non-positive Ricci curvature corollary}


\begin{corollary}\label{trivial}
Let $\left(M^{n},g,X,\lambda \right)$ be an $n$-dimensional complete, simply connected  
Ricci soliton of finite volume without boundary,
without conjugate points with the length of the potential field $X$, a positive constant. Let the scalar curvature $S$ be constant along the
integral curves of $X$ and the function $f=\left\langle \nabla _{v}\phi
X,v\right\rangle $ is integrable on the unit tangent bundle $SM^{n}$ of $M^{n}$. Then the Ricci curvature of $M$ along $X$ cannot be nonpositive
and consequently, in this case, the soliton must be shrinking.
\end{corollary}
\begin{proof}
By Theorem \ref{main}, the soliton is trivial, so $X$ is a Killing
field and $\operatorname{Ric}=\lambda g$; in particular
$\operatorname{Ric}(X,X)=\lambda\|X\|^2$. Set $h=\tfrac12\|X\|^2$.
Since $X$ is Killing,
\[
\Delta h = \|\nabla X\|^2 - \operatorname{Ric}(X,X)=0.
\]
If $\operatorname{Ric}(X,X)\le 0$, then $\|\nabla X\|^2\le 0$, so
$\nabla X=0$; hence $X$ is parallel and $\operatorname{Ric}(X,X)=0$.
By the hypotheses on $M$ (no conjugate points), every geodesic is a
line, so $M=N\times\mathbb{R}$, which has infinite volume, a
contradiction. Therefore $\operatorname{Ric}(X,X)>0$, and since
$\|X\|>0$ this gives $\lambda>0$; the soliton is shrinking.
\end{proof}

\bigskip
\begin{remark}
    From the Corollary \ref{trivial}, we conclude that 
$M$ cannot be a harmonic or asymptotically harmonic manifold, as they have strictly non-positive Ricci curvature.


\end{remark}

%%%%%%%%%%%%%%%%%%%%%%%%%%%%%%%%%%%%%%%%%%%%%%%%%%%$
\section{Applications to Almost Contact and Almost \texorpdfstring{$\alpha$}{alpha}-Cosymplectic Manifolds}
%%%%%%%%%%%%%%%%%%%%%%%%%%%%%%%%%%%%%%%%%%%%%%%%%%%%%


%-------------------------------------------------------%
\subsection{Reeb vector fields as soliton potentials}
%--------------------------------------------------------%

As the Reeb vector field $\xi $ of an almost contact metric manifold  $
\left( M^{2n+1},g, \varphi ,\xi ,\eta \right) $ is a unit vector field, as a
direct consequence of the above Theorem, we have:

\medskip 
Since $\|\xi\|=1$, one verifies that $\phi\xi=0$, whence the
integrability condition in Theorem~\ref{main} is automatically
satisfied.

\begin{corollary}\label{c1}
Let $(M^{2n+1},g,\xi,\lambda)$ be a non-compact Ricci soliton of finite
volume without boundary, where $\xi$ is the Reeb vector field of an
almost contact metric manifold $(M^{2n+1},g,\varphi,\xi,\eta)$. Suppose
that $S$ is constant along the integral curves of $\xi$ and that
\[
f=\langle\nabla_v\phi\xi,v\rangle
\]
is integrable on $SM^{2n+1}$, where $\phi$ is the skew-symmetric tensor
defined by \eqref{e1}. Then the soliton is trivial.
\end{corollary}


Now, observe that by Corollary \ref{c1}, if the $(2n+1)$-dimensional non-compact
Ricci soliton $\left( M^{2n+1},g,\xi ,\lambda \right) $ of finite volume,
where $\xi $ is the Reeb vector field of the almost contact metric manifold $%
\left( M^{2n+1},g, \varphi ,\xi ,\eta \right) $ such that the function $%
f=\left\langle \nabla _{v} \phi \xi ,v\right\rangle $ is integrable on the
unit tangent bundle $SM^{2n+1}$, then $\xi $ is Killing, that is $\left(
M^{2n+1},g, \varphi ,\xi ,\eta \right) $ is a $K$-contact Einstein manifold.
However, we cannot apply the result in \cite{boyer2001einstein} to make it a Sasakian Einstein owing
to non-compactness of $M^{2n+1}$. Yet, we have the following:

\medskip 

\begin{proposition}\label{prop:kcontact}
Let $(M^{2n+1},g,\varphi,\xi,\eta)$ be a non-compact almost contact
metric manifold of finite volume. Suppose that the Reeb vector field $\xi$
gives rise to a Ricci soliton $(M^{2n+1},g,\xi,\lambda)$. If the scalar
curvature $S$ is constant along the integral curves of $\xi$, then
$M^{2n+1}$ is a $K$-contact Einstein manifold.
\end{proposition}


%-------------------------------------------------------------------%
\subsection{Almost \texorpdfstring{$\alpha$}{alpha}-cosymplectic manifolds}
%-------------------------------------------------------------------%

\begin{lemma}\label{alpha-cosym}
Let $(M^n,g,\xi,\lambda)$ be an almost $\alpha$-cosymplectic
Ricci soliton with potential vector field $\xi$, the Reeb vector
field of the almost $\alpha$-cosymplectic manifold
$(M,g,\varphi,\xi,\eta)$ and scalar curvature $S$.
Then the soliton function satisfies
\[
n\lambda^2-(1+2S)\lambda+\|Q\|^2=0.
\]
Consequently,
\begin{equation}\label{lambdaeq}
\lambda=
\frac{(1+2S)\pm\sqrt{(1+2S)^2-4n\|Q\|^2}}{2n}.
\end{equation}
Moreover,
\[
1+2S>0
\qquad\text{and}\qquad
\lambda\ge0.
\]
Hence, the soliton is non-expanding (that is, shrinking or steady).
Furthermore, if $\lambda=0$, then $\xi$ is parallel and therefore
$M^3=\mathbb R^3/\Gamma_1$, while in general
$M^n=N^{n-1}\times(a,b)$, where $N$ has finite volume.
\end{lemma}

\begin{proof}
Note that $\eta$ is closed and therefore (\ref{e1})
and (\ref{e2}) imply that
\begin{eqnarray}\label{eq:nablaxi}
\nabla_{X} \xi = \lambda X - Q X \\ \nonumber
Q(\xi) = \lambda \xi.
\end{eqnarray}
Using the above equation, we compute
\begin{eqnarray}\label{curv}
R(U,V)\xi = - (\nabla Q)(U,V) + (\nabla Q)(V,U),
\end{eqnarray}
which, on contraction, yields
\begin{eqnarray}\label{e11}
\operatorname{Ric}(V,\xi)  = -\frac{1}{2} V(S) + V(S) =  \frac{1}{2} V(S)
\end{eqnarray}
that is,
\begin{eqnarray}\label{Qxi}
Q(\xi) = \frac{1}{2} \nabla S.
\end{eqnarray}
Note that \eqref{e5} together with the fact that
$\xi$ is closed, takes the form
\begin{eqnarray}\label{m1}
||Q||^2 - \frac{S^2}{n}  =
- \frac{1}{n} \left(S - n \lambda \right)^2 + \frac{1}{2}
\xi(S).
\end{eqnarray}
Now \eqref{eq:nablaxi} and \eqref{e11} respectively imply
$$\operatorname{Ric}(\xi, \xi) = \lambda \;\mbox{and} \;
\operatorname{Ric}(\xi, \xi) = \frac{1}{2} \xi(S).$$
Thus,
\begin{eqnarray}\label{m2}
\frac{1}{2} \xi(S) = \lambda.
\end{eqnarray}
Plugging (\ref{m2}) into (\ref{m1}) yields
\[
n\lambda^{2}-(1+2S)\lambda+\Vert Q\Vert^{2}=0 ,
\]
whose roots are given by \eqref{lambdaeq}.

It remains to show that $\lambda\geq0$. Since the quadratic equation admits the real root $\lambda$,
its discriminant must be non-negative at every point of $M^n$:
\[
(1+2S)^{2}\;\geq\;4n\Vert Q\Vert^{2}.
\]
On the other hand, the Cauchy--Schwarz inequality \eqref{e9} gives
$4n\Vert Q\Vert^{2}\geq 4S^{2}$. Combining the two inequalities,
\[
(1+2S)^{2}\geq 4S^{2},
\qquad\text{that is,}\qquad
1+4S\geq0 ,
\]
so that $S\geq-\frac14$ and hence
\[
1+2S\;\geq\;\frac12\;>\;0
\]
everywhere on $M^{n}$. By Vieta's formulas,
\[
r_1r_2=\frac{\|Q\|^2}{n}\ge0,
\qquad
r_1+r_2=\frac{1+2S}{n}>0.
\]
Since the roots are real and have a non-negative product,
they have the same sign (or one of them is zero). Because
their sum is positive, both roots must be non-negative.
As $\lambda$ is one of these roots, we conclude that
\[
\lambda\ge0.
\]
Hence, the soliton is non-expanding; that is, it is either shrinking or steady.

Finally, if $\lambda=0$, then the quadratic gives
$\Vert Q\Vert^{2}=0$, so $Q=0$ and, by \eqref{eq:nablaxi}, $\nabla\xi=0$.
Hence $\xi$ is parallel. Since $\xi$ is a non-vanishing parallel vector
field, the de Rham decomposition theorem implies that
\[
M^n=N^{n-1}\times(a,b),
\]
where $N$ has finite volume. In dimension three, this yields
\[
M^3=\mathbb{R}^3/\Gamma_1 .
\]
\end{proof}


\begin{corollary}\label{c:noexpanding}
There does not exist an expanding $n$-dimensional
almost $\alpha$-cosymplectic Ricci soliton 
$(M^n, g, \varphi, \xi, \eta)$.
\end{corollary}
\begin{proof}
Immediate from Lemma \ref{alpha-cosym}, which shows $\lambda\geq0$.
\end{proof}

It is worth noting that if an almost contact metric manifold $\left(
M^{2n+1},g, \varphi ,\xi ,\eta \right) $ is a $\alpha $-cosymplectic
manifold, then in this case the dual $1$-form $\eta $ to the Reeb vector
field $\xi $ is closed. Thus, for the Ricci soliton $\left( M^{2n+1},g,\xi
,\lambda \right) $, the operator $\phi $ appearing in equation (\ref{e1}) is zero
and accordingly the function $f=\left\langle \nabla _{v}\varphi \xi
,v\right\rangle =0$ is automatically integrable on the unit tangent bundle $%
SM^{2n+1}$ as required in Corollary \ref{c1}. Using this argument, we have



\begin{corollary}\label{c2}
Let $\left( M^3,g,\xi, \lambda \right)$ be a $3$-dimensional
non-compact Ricci soliton of finite volume without boundary with the potential 
field $\xi$, the Reeb vector field of the almost 
$\alpha$-cosymplectic manifold
 $\left( M^{3},g, \varphi ,\xi ,\eta \right) $. If the scalar curvature $S$ 
is constant along the integral curves of $\xi$, then $\left(M^{3},g,\xi,\lambda \right)$ is cosymplectic trivial and steady Ricci soliton. Consequently, $M^3$ is flat. If, in addition, $M^3$ is complete, then
$M^3=\mathbb{R}^3/\Gamma$, where $\Gamma$ is a Bieberbach group, that
is, a discrete torsion-free cocompact subgroup of the isometry group of
$\mathbb{R}^3$. In particular, such a complete finite-volume flat
manifold is compact.
\end{corollary}


\begin{proof}
By Corollary \ref{c1}, the Reeb vector field 
$\xi$ is Killing, which implies that $\nabla \xi$ is skew-symmetric. But for almost $\alpha$ cosymplectic
manifold $\nabla \xi$ is symmetric by \cite{li2022ricci}. Hence,
 $\nabla \xi = 0$, which implies that $\xi$ is parallel 
 vector field and, in turn, it is a Killing vector field. 
  The Levi-Civita connection on almost $\alpha$-cosymplectic $3$-manifold $M^{3}(\varphi,\xi,\eta,g)$ is described explicitly
 on p. 70, \cite{perrone2018classification}. We have
$$\nabla_e\xi =\alpha e-\sigma \varphi e = 0.$$
This implies that $\alpha = 0 = \sigma$.
Therefore, $M^3$ is cosymplectic and because we have
$$ Q\xi=-(2\alpha^2+\operatorname{tr}h^2)\xi+(2b\sigma-e(\sigma))\varphi e-(2a\sigma+(\varphi e)(\sigma))e,$$  
$\operatorname{Ric}(\xi, X) = 0$ for any $X$ and therefore,  $\lambda = 0$ and the soliton is  steady. \\
Substituting $\nabla\xi=0$ and $\lambda=0$ into the soliton equation
\eqref{e0}, we obtain
\begin{equation*}
    g(\nabla_{U}\xi,V)+g(U,QV)= \lambda \; g(U,V) = 0,
\end{equation*}
which simplifies to
\begin{equation*}
    g(U,QV)= 0.
\end{equation*}
Hence, $M^3$ is Ricci-flat and, in turn, flat. If $M^3$ is complete, then by the Bieberbach theorem
$M^3=\mathbb{R}^3/\Gamma$, where $\Gamma$ is a Bieberbach group.
Since a complete flat manifold of finite volume is compact, the quotient
is compact.
\end{proof}

\subsection{Transversal potential fields on almost \texorpdfstring{$\alpha$}{alpha}-cosymplectic \texorpdfstring{$3$}{}-manifolds}


Homogeneous almost $\alpha$-cosymplectic $3$-manifolds admitting Ricci solitons are studied in \cite{li2022ricci}.
In particular, the manifold in the aforementioned result must be homogeneous.


\begin{theorem}\label{tr1}
Let $\left(M^3, g, V, \lambda \right)$ be a $3$-dimensional non-compact almost  $\alpha$-cosymplectic 
Ricci soliton of finite volume without boundary with the potential
field $V$, a unit vector field, transversal to $\xi$, the Reeb vector field. Then $V$ is a trivial Ricci soliton if the scalar curvature $S$
is constant along the integral curves of $V$,
$a, b, \mu, \sigma$  are  bounded functions on $M$, and derivative of one of the coefficients of $V$ is bounded. In this case, $M$ is homogeneous and has constant curvature and $M^3 = {\mathbb R}^3/{\Gamma}_1$ or ${\mathbb H}^3/{\Gamma}_2$ is a compact steady trivial Ricci soliton or expanding trivial Ricci soliton. Here, ${\Gamma}_1$ and ${\Gamma}_2$ are, respectively, compact subgroups of the isometry group of $\mathbb{R}^3, \mathbb{H}^3$  acting transitively. 
 \end{theorem}
 
\begin{proof}
As $M^3$ is almost $\alpha$ cosymplectic manifold,
$\eta(X) = g(X, \xi)$ satisfies  $d \eta = 0$ and 
$d \Phi = 2 \alpha \eta \wedge \Phi$, where 
$\Phi(X,Y) = g(X, \varphi Y)$, where $\varphi$ is an almost contact structure. Therefore, this implies that 
\begin{equation}\label{l1}
d \Phi(E, e, F) = 2 \alpha \; \eta(E) \;  g(\varphi(e), F)
\end{equation}
\begin{equation}\label{l2}
d \Phi(E, \varphi e, F) = - 2 \alpha \; \eta(E) \; g(e, F) 
\; g(\varphi(e), F).
\end{equation}
As $V$ is a transversal unit vector field, we have 
$V = f_1 \; e + f_2 \; \varphi e$, where ${f_1}^2 + {f_2}^2 =1$.  Hence, clearly $f_1, f_2$ are bounded. Using (\ref{l1}) and (\ref{l2}) we obtain,  
\begin{equation}\label{l3}
d \Phi(E, V, F) = 2 \alpha \; \eta(E) \;  g(f_1 \varphi(e) - f_2 e, F).
\end{equation}
Therefore, we need to obtain conditions under which 
\begin{equation}\label{l4}
g\left({\nabla}_{v}(f_1 \varphi(e) - f_2 e), v \right)
\end{equation}
is integrable on $SM$. On expanding (\ref{l4}) we obtain,
\begin{equation}\label{l5}
g\left({\nabla}_{v}(f_1 \varphi(e) - f_2 e), v \right) =
({\nabla}_{v} f_1) a_2 + f_1  g({\nabla}_{v} \varphi e, v)
- ({\nabla}_{v} f_2) a_1 - f_2 g({\nabla}_{v} e, v),
\end{equation}
where $$v = a_1 e + a_2 \varphi e + a_3  \xi \; \mbox{with} \;
{a_1}^2 + {a_2}^2 + {a_3}^2 =1. $$ 
Now we compute $g(\nabla_{v} \varphi e, v)$ and 
$g(\nabla_{v}e, v)$ using connection coefficients as described in Theorem \ref{lc}, we obtain:
\begin{equation}\label{l6}
g({\nabla}_{v} \varphi e, v) = {a_1}^2 a + a_1 a_3 \sigma
+ a_1 a_2 b - a_2 a_3 \alpha - a_1 a_3 \mu.
\end{equation}
Using the connection coefficients, we obtain
\begin{equation}\label{l8}
g({\nabla}_{v} e, v) = - {a_2}^2 b  + a_2 a_3 \sigma - a_1 a_3 \alpha + a_3 a_2 \mu.
\end{equation}
Now we estimate the RHS of (\ref{l5}) using all the above equations. 
Note that by hypothesis, we have, 
${a_1}^2 + {a_2}^2 + {a_3}^2 = 1,$ consequently all
of $a_i$ are bounded. Since $f_1^2 + f_2^2 = 1$, if $df_1$ is bounded, then $df_2$ is also
bounded. Hence, the right-hand sides of (\ref{l6}) and (\ref{l8}) are
bounded by hypothesis, and consequently the right-hand side of (\ref{l5}) is bounded. Hence, RHS of (\ref{l6}),(\ref{l8}) are bounded by hypothesis and
consequently RHS of (\ref{l5}) is bounded.
And also, as $M$ is of finite volume, $SM$ is so.
Thus, the required expression (\ref{l4}) is integrable 
on $SM$. Hence, by Theorem \ref{main} $V$ is trivial 
Ricci soliton. This implies that 
$\operatorname{Ric}=\lambda g$ with $\lambda = \frac{r}{3}$.
But then $\lambda$ must be constant and hence $r$
is constant. By Proposition \ref{r} 
$$r=-6 \alpha^{2}-\operatorname{tr} h^{2}-2\left(a^{2}+b^{2}\right)-2(\varphi e)(a)+2 e(b),$$
hence $a, b, \sigma$ are all constants. Since $\alpha=0$ (forced by the cosymplectic structure) and
$a,b,\sigma$ are constants, the formula for $r$ in
Proposition~\ref{r} gives
\[
r=-\operatorname{tr}h^2-2(a^2+b^2)\le0.
\] If $r = 0$, then $M^3 = {\mathbb R^3}/{\Gamma_{1}}$,
and if $r < 0$, then $M^3 = {\mathbb H^3}/{\Gamma_{2}}$.\\
Thus, it is compact.
\end{proof}

From Corollary \ref{c2} and Theorem \ref{tr1} we obtain the complete characterisation of 
Ricci soliton on a non-compact, complete Riemannian manifold of finite volume.    

\begin{theorem}\label{tr2}
Let  $\left(M^3, g, V, \lambda \right)$ be a $3$-dimensional almost $\alpha$-cosymplectic non-compact Ricci soliton of finite volume without boundary with the potential field $V$, a unit vector field. Suppose that the scalar curvature $S$
is constant along the integral curves of $V$,
$a, b, \mu, \sigma$  are  bounded functions on $M$, and the derivatives of one of the coefficients of $V$ are bounded. In this case, $M$ is homogeneous of constant curvature  and $M^3 = {\mathbb R}^3/{\Gamma}_1$ or ${\mathbb H}^3/{\Gamma}_2$,  is compact steady trivial Ricci soliton.  
\end{theorem}

\begin{remark}
    Note that the expanding case in Theorem \ref{tr1} does not contradict
Corollary \ref{c:noexpanding}: the latter concerns the Reeb potential,
whereas here the potential $V$ is transversal to $\xi$.
\end{remark}
%%%%%%%%%%%%%%%%%%%%%%%%%%%%%%%%%%%%%%%%%%%%%%%%%%%$
\section{Contact Ricci solitons in dimension three}
%%%%%%%%%%%%%%%%%%%%%%%%%%%%%%%%%%%%%%%%%%%%%%%%%%%%%


\begin{theorem}\label{t:contact}
Let $(M^3, g, \xi, \lambda)$ be a $3$-dimensional contact Ricci
soliton with potential field $\xi$, the Reeb vector field of the
contact metric manifold $(M^3, g, \varphi, \xi, \eta)$. Then
$\lambda = 2$, the tensor $h$ vanishes, and the soliton is trivial.
Moreover, $M^3$ is a Sasakian Einstein manifold of constant sectional
curvature $1$. In particular, there exists no steady or expanding
$3$-dimensional contact Ricci soliton with potential the Reeb vector
field. If, in addition, $M^3$ is complete, then
$M^3=\mathbb{S}^3/\Gamma$ for a finite group $\Gamma$ acting freely and
isometrically, and hence $M^3$ is compact.
\end{theorem}
\begin{proof}
On the contact metric manifold $(M^3, g, \varphi, \xi, \eta)$ we have,
in the convention of \cite{blair2002riemannian},
$d\eta(U,V) = g(U, \varphi V)$, $U, V \in \chi(M^3)$, where the
exterior derivative carries the factor $\frac{1}{2}$. With respect to
the convention used in equation (\ref{e1}), namely
$d\eta(U,V) = U\eta(V) - V\eta(U) - \eta([U,V])$, the contact
condition reads $d\eta = 2\Phi$, and therefore the skew-symmetric
operator $\phi$ of equation (\ref{e1}) is identified as
$\phi = -\varphi$. (This is confirmed directly on $\mathbb{S}^3$,
where $\nabla_U \xi = -\varphi U$, $Q = 2I$ and $\lambda = 2$.)
Equation (\ref{e1}) now reads
\begin{eqnarray}\label{m3}
{\nabla}_{X} \xi = \lambda X - QX - \varphi X.
\end{eqnarray}
Thus we have,
\begin{eqnarray}\label{m4}
R(X,Y)\xi = - (\nabla Q)(X,Y) + (\nabla Q)(Y, X)
- (\nabla \varphi) (X,Y) + (\nabla \varphi) (Y,X),
\end{eqnarray}
which, on contraction, yields
\begin{eqnarray}\label{m5}
\operatorname{Ric}(Y, \xi) = \frac{1}{2}Y(S) + g\left(Y, {\sum}_{i}
(\nabla \varphi)(e_i, e_i)\right),
\end{eqnarray}
that is,
\begin{eqnarray}\label{m6}
Q(\xi) = \frac{1}{2} \nabla S + \sum_{i}(\nabla \varphi )(e_i, e_i).
\end{eqnarray}
Choose an adapted frame $\{e, \varphi e, \xi\}$; we have the
following structure equations:
\begin{equation}\label{m7}
    \begin{cases}
       & \nabla_ee= a \varphi e +b \xi,~~~~ \nabla_{\varphi e} e=c \varphi e+ d\xi,~~~~ \nabla_\xi e= f \varphi e + s \xi,  \\
       & \nabla_e \varphi e= -ce+\mu\xi,~~~~ \nabla_{\varphi e} \varphi e= -ce +\nu\xi,~~~~ \nabla_\xi \varphi e=-fe + h \xi ,  \\
       & \nabla_e\xi =- b e - \mu \varphi e,~~~~ \nabla_{\varphi e}\xi=-d e - \nu \varphi e ,~~~~ \nabla_\xi\xi= -se - h \varphi e,
    \end{cases}
  \end{equation}
where $a, b, c, \cdots, h$ are smooth functions. Using
(\ref{m7}) in (\ref{m3}), we arrive at
\begin{equation}\label{m8}
    \begin{cases}
       & Q(e) = (\lambda+b)e + (\mu - 1) \varphi e,   \\
       & Q(\varphi e) = (d + 1)e +(\lambda + \nu)\varphi e,  \\
       & Q(\xi) = se + h \varphi e + \lambda \xi.
    \end{cases}
\end{equation}
On using symmetry of $Q$ in equation (\ref{m8}), one
confirms $s=0$ and $h = 0$, and in particular we have
\begin{equation}\label{m9}
Q(\xi) = \lambda \xi
\end{equation}
and $\mu - 1 = d + 1$ gives
\begin{equation}\label{m10}
\mu - d = 2.
\end{equation}
Now using equation (\ref{m7}) with $s = h = 0$, we compute
\begin{eqnarray}\label{m11}
\sum_{i} (\nabla \varphi)(e_i, e_i) & = &
(\nabla \varphi)(e, e) + (\nabla \varphi)(\varphi e, \varphi e) + (\nabla \varphi)(\xi, \xi)\\ \label{sumphi}
& = & (\mu - d) \xi = 2\xi,
\end{eqnarray}
by virtue of (\ref{m10}). Hence, on plugging equations
(\ref{m6}), (\ref{m9}) and (\ref{sumphi}), we conclude
\begin{equation}\label{gradS}
\nabla S = 2(\lambda - 2) \xi.
\end{equation}
Taking the covariant derivative and using (\ref{m3}) we get
\begin{eqnarray}
H_S(X) = 2(\lambda - 2)(\lambda X - QX - \varphi X),
\end{eqnarray}
that is,
$$ \left(H_{S} - 2\lambda (\lambda - 2)I + 2(\lambda - 2) Q \right) X
= -2 (\lambda - 2) \varphi X.$$
The L.H.S. is symmetric, whereas the R.H.S. is skew-symmetric, which
implies that
$$ (\lambda - 2) \varphi X = 0, \quad X \in \chi(M^3).$$
Since $\varphi \neq 0$ on a contact manifold, we conclude
$\lambda = 2$.

Now, by (\ref{m9}), $\operatorname{Ric}(\xi,\xi) = \lambda = 2$. On a contact metric
$3$-manifold, one has
\[
\operatorname{Ric}(\xi,\xi)=2-\operatorname{tr}h^2
\]
\cite{blair2002riemannian}. Hence $\operatorname{tr}h^2=0$, that is,
$h=0$. Therefore $M^3$ is $K$-contact. In dimension three, every
$K$-contact manifold is Sasakian. Since $\xi$ is Killing, we have
$\mathcal{L}_\xi g=0$, and the soliton equation (\ref{e0}) gives
$\operatorname{Ric}=2g$. Hence, the soliton is trivial, and $M^3$ is Einstein. Since a
$3$-dimensional Einstein manifold has constant sectional curvature,
$M^3$ has constant sectional curvature $1$. If $M^3$ is complete, then
its universal cover is $\mathbb{S}^3(1)$, and therefore
$M^3=\mathbb{S}^3/\Gamma$ for a finite group $\Gamma$ acting freely and
isometrically. In particular, $M^3$ is compact.
\end{proof}





\bigskip



%%%%%%%%%%%%%%%%%%%%%%%%%%%%%%%%%%%%%%%%%%%%%%%%%%%$
\section{Examples and Sharpness}\label{sec:examples}
%%%%%%%%%%%%%%%%%%%%%%%%%%%%%%%%%%%%%%%%%%%%%%%%%%%%%




In this section, we construct several examples which show that the
hypotheses appearing in our main results are essential, and which
connect the present work with the classification of Ricci solitons on
homogeneous almost $\alpha$-cosymplectic three-manifolds obtained by
Li and Liu \cite{li2022ricci}.

\begin{remark}[Notation for {\cite{li2022ricci}}]\label{rem:notation}
When quoting \cite{li2022ricci}, the following dictionary must be kept
in mind: the soliton constant denoted $S$ in \cite{li2022ricci} is our
$\lambda$; the eigenvalue of the tensor $h$ denoted $\lambda$ in
\cite{li2022ricci} is our $\sigma$; the scalar curvature denoted $r$
in \cite{li2022ricci} is our $S$. The $\varphi$-basis
$\{\xi, e, \varphi e\}$ and the structure functions $a, b, \mu$
coincide with ours, since both papers use the frame of Perrone
\cite{perrone2018classification}.
\end{remark}

%---------------------------------------------------------------------
\subsection{A warped product non-example: the finite-volume Kenmotsu cusp}
%---------------------------------------------------------------------

\begin{example}\label{ex:cusp}
Let $M^{3}=(-\infty,0)\times\mathbb{T}^{2}$, where $\mathbb{T}^{2}$ is
the flat two-torus with coordinates $(x,y)$, equipped with the warped
product metric
\begin{equation}\label{eq:cuspmetric}
g \;=\; ds^{2}+e^{2s}\left(dx^{2}+dy^{2}\right).
\end{equation}
Define the almost contact metric structure by
\[
\xi=\partial_{s},\qquad
\eta=ds,\qquad
\varphi(\partial_{x})=\partial_{y},\qquad
\varphi(\partial_{y})=-\partial_{x},\qquad
\varphi(\xi)=0 ,
\]
where $\varphi$ acts as the rotation by $\pi/2$ on $\ker\eta$. Since
$\eta=ds$ we have $d\eta=0$, and the fundamental $2$-form
$\Phi=-e^{2s}\,dx\wedge dy$ satisfies
\[
d\Phi \;=\; -2\,ds\wedge e^{2s}\,dx\wedge dy \;=\; 2\,\eta\wedge\Phi .
\]
Hence $\left(M^{3},g,\varphi,\xi,\eta\right)$ is an almost
$\alpha$-cosymplectic manifold with $\alpha=1$; in fact it is Kenmotsu
and $h=0$. Moreover:
\begin{enumerate}
\item[(i)] $M^{3}$ is non-compact, without boundary, and of
\emph{finite volume}:
\[
\operatorname{Vol}\left(M^{3}\right)
=\operatorname{Vol}\left(\mathbb{T}^{2}\right)
\int_{-\infty}^{0}e^{2s}\,ds
=\frac{1}{2}\operatorname{Vol}\left(\mathbb{T}^{2}\right)<\infty .
\]
\item[(ii)] The metric \eqref{eq:cuspmetric} has constant sectional
curvature $-1$, so that $\operatorname{Ric}=-2g$ and $Q\xi=-2\xi$.
\item[(iii)] $\nabla_{U}\xi=U-\eta(U)\xi$ for all
$U\in\Gamma\left(TM^{3}\right)$, whence
$\frac{1}{2}\mathcal{L}_{\xi}g=g-\eta\otimes\eta$ and
\[
\frac{1}{2}\mathcal{L}_{\xi}g+\operatorname{Ric}\;=\; -g-\eta\otimes\eta .
\]
\end{enumerate}
Evaluating the last identity on $(\xi,\xi)$ gives $\lambda=-2$,
whereas evaluating it on horizontal unit vectors gives $\lambda=-1$.
Consequently, there exists \emph{no} constant $\lambda$ for which
$\left(M^{3},g,\xi,\lambda\right)$ is a Ricci soliton: the unit Reeb
field of the cusp does not generate a Ricci soliton. This is in exact
agreement with Lemma \ref{alpha-cosym} and Corollary
\ref{c:noexpanding}, since $Q\xi=-2\xi$ would force the value
$\lambda=-2<0$, contradicting the non-negativity of the soliton
function in equation \eqref{lambdaeq}.

This example shows, on the one hand, that the class of non-compact
almost $\alpha$-cosymplectic manifolds of finite volume without
boundary is non-vacuous and, on the other hand, that the Ricci soliton
condition in our results is a genuine geometric restriction on this
class.
\end{example}

%---------------------------------------------------------------------
\subsection{An almost \texorpdfstring{$\alpha$}{alpha}-cosymplectic
example attaining Corollary \ref{c2}}
%---------------------------------------------------------------------

\begin{example}\label{ex:torus}
By Theorem 3.3 of \cite{li2022ricci}, the abelian Lie group
$\mathbb{R}^{3}$ (as well as the group $\widetilde{E}^{2}$), equipped
with a flat left-invariant \emph{cosymplectic} structure, carries a
contact Ricci soliton $\left(\mathbb{R}^{3},g,\xi,0\right)$: the Reeb
field $\xi$ is parallel, hence a unit Killing field, $\operatorname{Ric}=0$, and the
soliton is steady and trivial, in agreement with Corollary 3.4 of
\cite{li2022ricci} and with our Corollary \ref{c2}. Passing to the
quotient by the integer lattice $\mathbb{Z}^{3}$, whose translations
preserve the left-invariant structure, we obtain the flat torus
\[
\mathbb{T}^{3}=\mathbb{R}^{3}/\mathbb{Z}^{3},
\]
a finite-volume cosymplectic manifold carrying the steady trivial
contact Ricci soliton and realising precisely the conclusion
$M^{3}=\mathbb{R}^{3}/\Gamma$ of Corollary \ref{c2}. All hypotheses of
Corollary \ref{c2} are satisfied: $\left\Vert\xi\right\Vert=1$, the
scalar curvature $S=0$ is constant, and since $d\eta=0$ the operator
$\phi$ of equation (\ref{e1}) vanishes, so that the integrability
condition on the unit tangent bundle holds automatically.
\end{example}

\begin{remark}\label{rem:bieberbach}
By the Bieberbach theorem, a complete flat manifold of finite volume
is necessarily compact. Hence, Corollary \ref{c2}, and likewise
Theorems \ref{tr1} and \ref{tr2}, may be equivalently restated as
\emph{non-existence} results: there is no genuinely non-compact Ricci
soliton of finite volume in this class, and the conclusion is attained
only on compact quotients such as the torus $\mathbb{T}^{3}$ of
Example \ref{ex:torus}.
\end{remark}

%---------------------------------------------------------------------
\subsection{Necessity of the finite-volume hypothesis}
%---------------------------------------------------------------------

\begin{example}\label{ex:sol3}
The finite-volume hypothesis cannot be removed from our results.
Following Theorem 4.1(i) of \cite{li2022ricci}, consider the
unimodular Lie group
\[
Sol_{3}=\mathbb{R}^{2}\rtimes_{A}\mathbb{R},
\qquad
A=\begin{pmatrix} 0 & -\sigma\\ -\sigma & 0\end{pmatrix},
\]
equipped with its standard left-invariant \emph{strictly almost
cosymplectic} structure, that is, $\alpha=0$, $h\neq 0$, $\sigma$ a
non-zero constant, and $a=b=\mu=0$. By \cite{li2022ricci}, $Sol_{3}$
carries a Ricci soliton $\left(Sol_{3},g,V,\lambda\right)$ whose
potential field $V=f_{1}e+f_{2}\varphi e$ is \emph{transversal} to
$\xi$, with
\[
\lambda=-2\sigma^{2}<0,
\]
that is, an expanding, non-trivial Ricci soliton; indeed $Sol_{3}$ is
homogeneous but not Einstein. Here $M$ is simply connected,
non-compact, and of \emph{infinite volume}; the scalar curvature
$S=-2\sigma^{2}$ is constant, in particular constant along the
integral curves of $V$; and $a,b,\mu,\sigma$ are constants. Thus,
every hypothesis of Theorems \ref{tr1} and \ref{tr2} is satisfied
\emph{except} finite volume, while the conclusion (triviality and
compactness) fails. Hence, the finite-volume hypothesis is essential.
\end{example}

\begin{remark}[Consistency with Theorem \ref{tr1}]\label{rem:consistency}
It is instructive to observe that the potential field $V$ of Example
\ref{ex:sol3} cannot have constant length, exactly as Theorem
\ref{tr1} predicts. Indeed, specialising equations (4.3)--(4.7) of
\cite{li2022ricci} to $Sol_{3}$ (where $\alpha=a=b=\mu=0$) yields
\[
e(f_{1})=-2\sigma^{2},\qquad
\varphi e(f_{2})=-2\sigma^{2},\qquad
\varphi e(f_{1})+e(f_{2})=0 .
\]
Suppose $f_{1}^{2}+f_{2}^{2}$ were constant. Differentiating along $e$
and $\varphi e$ gives
\[
-4\sigma^{2}f_{1}+2f_{2}\,e(f_{2})=0,
\qquad
2f_{1}\,\varphi e(f_{1})-4\sigma^{2}f_{2}=0 .
\]
Substituting $\varphi e(f_{1})=-e(f_{2})$ into the second identity and
eliminating $e(f_{2})$ leads to $f_{1}^{2}+f_{2}^{2}=0$, that is,
$V=0$, which would force $Sol_{3}$ to be Einstein --- a contradiction.
Therefore $\left\Vert V\right\Vert$ is necessarily non-constant on
$Sol_{3}$: the solitons of \cite{li2022ricci} lie strictly outside the
constant-length regime governed by Theorem \ref{main}, which is why no
contradiction with our results arises.
\end{remark}

%---------------------------------------------------------------------
\subsection{Sharpness of Theorem \ref{t:contact}}
%---------------------------------------------------------------------

\begin{example}\label{ex:sphere}
Consider the unit sphere $\mathbb{S}^{3}(1)$ with its standard
Sasakian structure $\left(\varphi,\xi,\eta,g\right)$ arising from the
Hopf fibration. Then $\xi$ is a unit Killing field and $\operatorname{Ric}=2g$, so
that
\[
\frac{1}{2}\mathcal{L}_{\xi}g+\operatorname{Ric} = 2g,
\]
that is, $\left(\mathbb{S}^{3},g,\xi,2\right)$ is a shrinking contact
Ricci soliton with potential field the Reeb vector field; compact
quotients $\mathbb{S}^{3}/\Gamma$ furnish further examples. This shows
that the shrinking regime in Theorem \ref{t:contact} is genuinely
occupied and that a non-existence statement for contact Ricci solitons
with Reeb potential can only concern the \emph{non-trivial} ones. On
$\mathbb{S}^{3}$ one checks directly that $\nabla_{U}\xi=-\varphi U$,
so that comparison with equation (\ref{e1}) (where $Q=2I$ and
$\lambda=2$) identifies the skew-symmetric operator as
$\phi=-\varphi$, and the structure functions of the adapted frame
\eqref{m7} satisfy $\mu=1$, $d=-1$.
\end{example}

\begin{remark}\label{rem:sharpness}
The conclusion of Theorem \ref{t:contact} is sharp in the following
sense. Once the potential field is allowed to be \emph{transversal} to
the Reeb field and the finite-volume condition is dropped, expanding
examples do exist: by Theorem 4.4 of \cite{li2022ricci}, the
hyperbolic space $\mathbb{H}^{3}(-\alpha^{2})$ carries an expanding
homogeneous almost $\alpha$-Kenmotsu structure with a transversal
potential vector field, and by Theorem 4.1(i) of \cite{li2022ricci}
the group $Sol_{3}$ of Example \ref{ex:sol3} carries an expanding
strictly almost cosymplectic one. Together with Example
\ref{ex:sphere} (shrinking, Reeb potential, compact) and Example
\ref{ex:torus} (steady, Reeb potential, flat), all three regimes ---
shrinking, steady, expanding --- are realised, each precisely in the
geometric situation permitted by our results.
\end{remark}

%%%%%%%%%%%%%%%%%%%%%%%%%%%%%%%%%%%%%%%%%%%%%%%%%%%%%%
\section*{Concluding Remarks}
%%%%%%%%%%%%%%%%%%%%%%%%%%%%%%%%%%%%%%%%%%%%%%%%%%%%%

The results obtained above show that the finite-volume assumption imposes
strong rigidity on Ricci solitons whose potential field has constant
length. In particular, in the almost contact setting, Reeb-potential
Ricci solitons are forced to be trivial under natural scalar curvature
and integrability assumptions. The examples in Section \ref{sec:examples} show that the finite-volume hypothesis and the constant-length condition
are essential.

It would be interesting to investigate whether the integrability
assumption on $\operatorname{div}(\phi X)$ can be weakened, or whether similar rigidity phenomena hold for Ricci almost solitons and other generalised soliton structures.

\bibliographystyle{alpha}
\bibliography{ref}


\end{document}
